% TOP_PROBLEM: {"problemId": "problem.provisional-top500-omission-wave-004.004", "canonicalTitle": "Original Dubrovin conjecture for Fano manifolds", "releaseRank": 316, "primaryDomain": "geometry_topology", "primaryDomainLabel": "Geometry & topology", "displayStatus": "open", "releaseStatus": "open"}
% !TeX program = lualatex
% ATTEMPT_STATUS: unresolved
\documentclass[11pt]{article}
\usepackage[margin=25mm]{geometry}
\usepackage{fontspec}
\setmainfont{Latin Modern Roman}
\usepackage{amsmath,amssymb,mathtools}
\usepackage{unicode-math}
\setmathfont{Latin Modern Math}
\usepackage{xurl,hyperref}
\hypersetup{colorlinks=true,urlcolor=blue}
\setcounter{secnumdepth}{0}
\setlength{\parindent}{0pt}
\setlength{\parskip}{5pt}
\setlength{\emergencystretch}{3em}
\begin{document}
\section*{\texorpdfstring{Original Dubrovin conjecture for Fano manifolds}{Original Dubrovin conjecture for Fano manifolds}}\label{316}
\subsection{Definitions and mathematical statement}
For a smooth projective complex Fano manifold $X$, the anticanonical line bundle is ample. Big quantum cohomology is the deformation of $H^*(X)$ whose product is defined by genus-zero Gromov--Witten invariants; generic semisimplicity means its generic algebra is a product of fields. In $D^b\operatorname{Coh}(X)$, an exceptional collection $(E_1,\ldots,E_N)$ has $\operatorname{Ext}^*(E_i,E_i)={\mathbb{C}}$ in degree zero and $\operatorname{Ext}^*(E_j,E_i)=0$ for $j>i$; it is full if it generates the category. The original Dubrovin package first asserts generic semisimplicity exactly when a full exceptional collection exists. Under the analytic and admissible-phase hypotheses of the cited sections, it further requires a collection with
\[
 S_{ij}=\chi(E_i,E_j),\qquad C=C'C'',
\]
where $S,C$ are Stokes and central connection matrices, $C''$ has the modified Chern-character columns of that source, and $C'$ commutes with cup product by $c_1(X)$. These normalization and convergence conditions are retained by reference; a modern Gamma-conjecture replacement is not substituted.
\par\smallskip\textit{Formulation and concept references:} \href{https://arxiv.org/pdf/1404.6407v4}{[S1]}.

\subsection{Short English statement}
Does the original Dubrovin package connect semisimple quantum cohomology, full exceptional collections, and the specified Stokes and connection matrices for every Fano manifold?
\subsection{Research attempt: projective-space algebra and the analytic matching gap}
\paragraph{1. Keep all parts of the original package visible.}
There are three logically separate requests: generic semisimplicity if and only if a full exceptional collection exists; identification of a Stokes matrix with its Euler matrix; and the specified factorization of the central connection matrix. Section 4.6 of [S1], checked in the HTML rendering [E1] on 27 September 2026, lists those original assertions. Its Proposition 4.6.6 states that Gamma Conjecture II implies the second and third parts under its hypotheses; it does not by itself prove the first equivalence for every Fano manifold. Theorem 5.0.1 and Corollary 5.0.2 treat projective spaces. Those analytic results are attributed to the source. The independent calculations below reproduce algebraic consistency checks and show why they do not determine Stokes data by themselves.

\paragraph{2. The Euler matrix gives a necessary integral structure.}
Let $(E_0,\ldots,E_r)$ be an exceptional collection. By exceptionality and semiorthogonality, its Euler matrix
\[
 G_{ij}=\chi(E_i,E_j)=\sum_k(-1)^k\dim\operatorname{Ext}^k(E_i,E_j)
\]
is upper triangular with every diagonal entry one. Hence $\det G=1$. Its classes in $K_0(X)$ are linearly independent: if $\sum_j a_j[E_j]=0$, pairing with each $E_i$ gives $Ga=0$, and invertibility forces $a=0$. When the collection is full, standard exceptional projection functors give a semiorthogonal decomposition into the categories generated by its individual objects, so $K_0(X)$ is freely generated by these classes. The fullness-to-generation statement uses that categorical decomposition; it is not inferred from an arbitrary collection of independent numerical classes.

Thus any proposed Stokes identification has an integral, unimodular upper-triangular matrix as its categorical target, in the appropriate ordered basis. Finding some matrix with these properties is only a necessary test. Many such matrices are not known to arise from exceptional collections, and a numerical Euler pairing does not identify the analytic asymptotic basis.

\paragraph{3. Generic semisimplicity for projective space computed explicitly.}
For $X=\mathbb P^n$, the standard small quantum-cohomology presentation is
\[
 QH^*(\mathbb P^n)\big|_q\cong\mathbb C[h]/(h^{n+1}-q).
\]
This enumerative presentation is a known input; the splitting calculation is elementary. For $q\ne0$, choose the $n+1$ distinct roots $r_j$ of $q$. The derivative $(n+1)h^n$ has no common root with $h^{n+1}-q$, so the quotient algebra is a product of $n+1$ copies of $\mathbb C$. Explicit primitive idempotents are
\[
 e_j=\frac1{n+1}\sum_{k=0}^{n}\left(\frac h{r_j}\right)^k.
\]
Evaluating at $h=r_\ell$ gives one if $\ell=j$ and zero otherwise, by the geometric-series formula. Therefore $e_je_\ell=\delta_{j\ell}e_j$ and $\sum e_j=1$.

At $q=0$, the algebra has the nonzero nilpotent $h$ and is not semisimple for $n\geq1$. This does not contradict generic semisimplicity. Conversely, finding nonsemisimplicity of a small quantum specialization for another Fano does not disprove generic semisimplicity of its big quantum cohomology. A special locus may be contained in the discriminant. For projective space, the displayed semisimple point already proves generic semisimplicity in the convergent family because nonvanishing of the algebra discriminant is an open condition.

For a parity-preserving supercommutative specialization on full cohomology, any odd element $a$ satisfies $a^2=-a^2$, hence $a^2=0$ over $\mathbb C$. A product of fields has no nonzero nilpotents, so such a semisimple full algebra cannot contain odd cohomology. This is a necessary parity obstruction, not a sufficient condition. It is stated for full cohomology; restricting by convention to the even part changes what this argument can test.

\paragraph{4. The categorical side for projective space.}
Consider the sequence $\mathcal O,\mathcal O(1),\ldots,\mathcal O(n)$. For $i,j$ in this range,
\[
 \operatorname{Ext}^k(\mathcal O(i),\mathcal O(j))
 =H^k(\mathbb P^n,\mathcal O(j-i)).
\]
The standard cohomology of line bundles says that $H^k(\mathcal O)=0$ for $k>0$, and every cohomology group of $\mathcal O(-a)$ vanishes for $1\leq a\leq n$. These facts prove exceptionality and the required vanishing when $j<i$. For $j\geq i$, only degree-zero cohomology contributes, giving
\[
 G_{ij}={n+j-i\choose n},\qquad G_{ij}=0\quad(j<i).
\]
The collection is full by Beilinson's resolution of the diagonal. One can see the mechanism from the exact complex on $\mathbb P^n\times\mathbb P^n$
\[
 0\to\mathcal O(-n)\boxtimes\Omega^n(n)\to\cdots
 \to\mathcal O(-1)\boxtimes\Omega^1(1)
 \to\mathcal O\boxtimes\mathcal O\to\mathcal O_\Delta\to0.
\]
It is the Koszul resolution of the diagonal, a regular zero locus of rank $n$. Applying the integral transform with this kernel to any bounded coherent complex expresses it through the line bundles $\mathcal O(-n),\ldots,\mathcal O$. Twisting the whole collection by $\mathcal O(n)$ yields the displayed sequence. This explains why fullness is stronger than matching the number of objects to a Betti number: the diagonal resolution provides generation for every object.

Thus the first part of the package holds for projective space by explicit known structures. For $\mathbb P^1$, the Euler matrix is
\[
 G=\begin{pmatrix}1&2\\0&1\end{pmatrix}.
\]
For $\mathbb P^2$ it is $\left(\begin{smallmatrix}1&3&6\\0&1&3\\0&0&1\end{smallmatrix}\right)$. These matrices give concrete targets for a Stokes calculation, not a derivation of that calculation.

\paragraph{5. The commutation condition on the central matrix is substantial.}
Let $H$ denote classical cup product by the hyperplane class on $H^*(\mathbb P^n)$, in the basis $1,h,\ldots,h^n$. It is a single nilpotent Jordan block. An operator $A$ commuting with $H$ is determined by $A(1)$, since
\[
 A(h^k)=A(H^k1)=H^kA(1).
\]
Writing $A(1)=\sum_{j=0}^n a_jh^j$ proves
\[
 A=a_0I+a_1H+\cdots+a_nH^n.
\]
Conversely every such polynomial commutes with $H$. It is invertible exactly when $a_0\ne0$, by triangularity. Since $c_1(\mathbb P^n)=(n+1)h$, commuting with cup product by $c_1$ is equivalent to commuting with $H$. This is an $(n+1)$-dimensional family of operators, far smaller than the unrestricted space of all matrices.

Use the modified Chern character of [S1], which multiplies the degree-$2k$ component of the ordinary Chern character by $(2\pi i)^k$. On $\mathbb P^1$, the columns for $\mathcal O,\mathcal O(1)$ form
\[
 C''=\begin{pmatrix}1&1\\0&2\pi i\end{pmatrix}.
\]
The commutant calculation gives $C'=\left(\begin{smallmatrix}a&0\\b&a\end{smallmatrix}\right)$ with $a\ne0$. Therefore the requested factorization, in these fixed conventions and this ordered collection, would have the form
\[
 C=C'C''=\begin{pmatrix}a&a\\b&b+2\pi i a\end{pmatrix}.
\]
The first-row equality and the specified second-row difference are genuine constraints. They are not automatic consequences of invertibility of $C$. Changing the collection, order, signs, or sector changes how the comparison is expressed; one cannot use an inconsistent normalization to claim a contradiction.

\paragraph{6. Why algebra eigenvectors do not determine Stokes matrices.}
The idempotents diagonalize quantum multiplication at a semisimple point. The Stokes matrix compares asymptotic fundamental solutions of a differential equation across sectors near an irregular singularity. It depends on analytic continuation and the choice of admissible phase. An eigenbasis of the leading coefficient gives formal exponential directions but does not determine the exponentially small changes between sectorial solutions. Solving the finite-dimensional algebra in Section 3 therefore leaves an analytic problem.

There is a matching categorical ambiguity. Mutating an exceptional pair changes its classes by an integral transformation: using the convention that the defining triangle is $\operatorname{RHom}(E,F)\otimes E\to F\to L_EF$,
\[
 [L_EF]=[F]-\chi(E,F)[E].
\]
This follows immediately by taking classes in the triangle defining $L_EF$. Under a basis change $Q$, the Euler matrix transforms as $Q^TGQ$. A Stokes matrix computed in one phase need not match the original unmutated line-bundle order; the source's projective-space theorem expressly allows mutations. Matching one ad hoc matrix entry after an arbitrary basis change does not establish the whole package.

The projective-space analytic conclusion is known from [S1], Corollary 5.0.2. This attempt does not rederive its asymptotic analysis. Its calculations independently verify the algebraic and normalization side against which that analysis must be compared.

\paragraph{7. Diagnostics and remaining gap.}
The saved exact diagnostic checks the binomial Euler matrices and determinant one, the commutant dimension by rational linear algebra for small $n$, and the two idempotent identities for $\mathbb P^1$ at $q=1$. It does not compute Gromov--Witten invariants of arbitrary Fanos, certify convergence of big quantum cohomology, or integrate an irregular connection.

\textbf{UNRESOLVED.} The projective-space subcase is explained with attributed analytic input and explicit independent algebra. For a general Fano, neither implication in the first equivalence is proved here, and even obtaining both structures would leave the analytic identification of compatible Stokes and central matrices. Gamma Conjecture II is a conditional route to part of that task, not a replacement of the original requested package.

\subsection{Sources}
\begin{itemize}
\item[S1]S. Galkin, V. Golyshev and H. Iritani, Gamma classes and quantum cohomology of Fano manifolds: Gamma conjectures, arXiv:1404.6407v4, Section 4.6.
\url{https://arxiv.org/pdf/1404.6407v4}.
\textit{Formulation source.}
\item[E1]S. Galkin, V. Golyshev and H. Iritani, Gamma classes and quantum cohomology of Fano manifolds: Gamma conjectures, arXiv:1404.6407v4, HTML rendering, Sections 4.4--4.6 and 5.
\url{https://arxiv.org/html/1404.6407v4}.
\textit{Read 27 September 2026 for the exact original package, Proposition 4.6.6 and the projective-space Theorem 5.0.1/Corollary 5.0.2. General analytic hypotheses are retained.}
\end{itemize}
\end{document}
